% Generated draft: requires mathematical and source review.
\documentclass[11pt]{article}
\usepackage[margin=25mm]{geometry}
\usepackage{fontspec}
\setmainfont{Latin Modern Roman}
\usepackage{amsmath,amssymb,mathtools,unicode-math}
\setmathfont{Latin Modern Math}
\usepackage{enumitem,needspace,xurl,hyperref}
\setlength{\parindent}{0pt}
\setlength{\parskip}{5pt}
\setlength{\emergencystretch}{3em}
\newcommand{\N}{\mathbb{N}}
\newcommand{\Z}{\mathbb{Z}}
\newcommand{\Q}{\mathbb{Q}}
\newcommand{\R}{\mathbb{R}}
\newcommand{\C}{\mathbb{C}}
\newcommand{\F}{\mathbb{F}}
\newcommand{\A}{\mathbb{A}}
\newcommand{\PP}{\mathbb{P}}
\newcommand{\E}{\mathbb{E}}
\newcommand{\eps}{\varepsilon}
\newcommand{\dd}{\,\mathrm{d}}
\begin{document}
\section*{UnsolvedMath ID: 30005959}
\textbf{Source problem code:} catalogue entry
\subsection{Definitions and mathematical statement}
Let \(X\) be an irreducible holomorphic symplectic manifold: a compact, simply connected K\"ahler manifold with \(H^0(X,\Omega_X^2)=\C\sigma\), where \(\sigma\) is everywhere nondegenerate. Write \(\dim_{\C}X=2n\). Let \(\pi:X\to B\) be a Lagrangian fibration, meaning a proper surjective holomorphic map with connected fibres whose smooth fibres have complex dimension \(n\) and on which \(\sigma\) restricts to zero.

A Shafarevich--Tate twist \(\pi^{\phi}:X^{\phi}\to B\) is obtained by gluing local copies of this fibration using the fibrewise automorphisms admitted by the Shafarevich--Tate construction. Denote its group of twist parameters by \(\mathcal T_{\pi}\). This notation is a placeholder: the supplied record does not define the automorphism sheaf, parameter group, or equivalence relation used in that construction. The notation rendered as \(X\) and \(X_0\) in the extracted parameter-group expressions is denoted here by \(\mathcal T_{\pi}\) and \(\mathcal T_{\pi,0}\), respectively, to distinguish it from the manifold \(X\). The meaning of the subgroup \(\mathcal T_{\pi,0}\) must be recovered from the source.

A complex manifold is K\"ahler if it admits a closed real \((1,1)\)-form that is positive definite. The requested problem is: for each fibration \(\pi\) as above, give necessary and sufficient conditions on every admissible class \(\phi\in\mathcal T_{\pi}\) for \(X^{\phi}\) to be K\"ahler; equivalently, characterize
\[
\mathcal K_{\pi}=\{\phi\in\mathcal T_{\pi}:X^{\phi}\text{ is K\"ahler}\}.
\]

The supplied excerpt reports the following partial results, subject to the source's definitions and any hypotheses omitted from the extraction. If there exists \(r\in\Z_{>0}\) such that \(r\phi\in\mathcal T_{\pi,0}\), then \(X^{\phi}\) is K\"ahler. It also uses the Leray differential
\[
d_2^{0,2}:H^0(B,R^2\pi_*\Q_X)\longrightarrow H^2(B,R^1\pi_*\Q_X)
\]
and a supplied identification
\[
H^2(B,R^1\pi_*\Q_X)\simeq(\mathcal T_{\pi}/\mathcal T_{\pi,0})\otimes_{\Z}\Q.
\]
Let \(\bar\phi\) denote the image of \(\phi\) under this identification. Writing \(b_2(Y)=\dim_{\Q}H^2(Y,\Q)\), the excerpt reports:
\begin{itemize}
\item If \(\bar\phi\in\operatorname{im}(d_2^{0,2})\), then \(b_2(X^{\phi})=b_2(X)\), and there is a degree-two cohomology class on \(X^{\phi}\) whose restriction to a smooth fibre is ample.
\item If \(\bar\phi\notin\operatorname{im}(d_2^{0,2})\), then \(b_2(X^{\phi})=b_2(X)-1\), all degree-two cohomology classes on \(X^{\phi}\) restrict trivially to a smooth fibre, and \(X^{\phi}\) is not K\"ahler.
\end{itemize}
The coefficients for the degree-two classes in these two assertions, and the precise meaning of an ample restriction, are not specified in the excerpt. These assertions are preserved as supplied source claims, not independently verified theorems. In particular, the excerpt does not assert that membership of \(\bar\phi\) in \(\operatorname{im}(d_2^{0,2})\) is sufficient for K\"ahlerness. The missing construction and hypotheses prevent this draft from fixing a fully specified parameter domain.
\subsection{Short English statement}
For a fixed Lagrangian fibration on an irreducible holomorphic symplectic manifold, characterize exactly which Shafarevich--Tate twists have K\"ahler total spaces. The supplied source reports sufficient and obstructing conditions, but its extracted text omits definitions needed to specify the twist parameters precisely. Its partially solved status and literature assessment are supplied dated curation.
\subsection{Sources}
\begin{itemize}
\item Local UnsolvedMath registry, supplied record with permanent ID \(30005959\), title \emph{K\"ahler Shafarevich--Tate Twists}, exact source problem code \texttt{catalogue entry}, in the Oberwolfach Reports Open Problems collection. The registry gives proposed year 2024 and names no proposer.
\item Supplied report citation: \emph{Algebraic Geometry: Wall Crossing and Moduli Spaces, Varieties and Derived Categories} (2025), Oberwolfach Reports, pp.~1787--1854, DOI \texttt{10.4171/owr/2024/32}. Original report and supplied literature-source link: \url{https://doi.org/10.4171/owr/2024/32}. The extracted passage contains Question~1 and Theorems~2 and~3. Its first-person attribution to upcoming work has no named speaker in the supplied record.
\item Supplied literature source labelled ``Shafarevich--Tate twists primary paper'': \url{https://link.springer.com/article/10.1007/s00209-025-03776-w}. No authors or full bibliographic title are supplied for this link.
\item Literature assessment dated exactly \texttt{2026-08-22}: the supplied curation labels the problem partially solved and reports K\"ahlerness for broad torsion and related classes of twists while leaving the infinite-order component question open. This is supplied dated curation, not an independently verified literature assessment. The registry also labels the statement \texttt{corrected\_verified} and reports reconstruction and checking against the official report; those are supplied provenance claims.
\item References attached to the apparently unrelated syzygetic-case fragment in the extracted text: [1] C.~Brooke, S.~Frei, L.~Marquand, X.~Qin, \emph{Birational geometry of Fano varieties of lines on cubic fourfolds containing pairs of cubic scrolls} (2024), \url{https://arxiv.org/abs/2407.18904}; [2] B.~Hassett, Y.~Tschinkel, \emph{Flops on holomorphic symplectic fourfolds and determinantal cubic hypersurfaces}, J.~Inst.~Math.~Jussieu \textbf{9} (2010), 125--153; [3] L.~Marquand, \emph{Cubic fourfolds with an involution}, Trans.~Amer.~Math.~Soc. \textbf{376} (2023), 1373--1406. Their relevance to this twist question is not established by the supplied record.
\end{itemize}
\end{document}
